% Part: first-order-logic
% Chapter: model-theory
% Section: theory-of-m

\documentclass[../../../include/open-logic-section]{subfiles}

\begin{document}

\olfileid{mod}{bas}{thm}

\section{اک \printtoken{S}{structure} دا نظریہ}

ہر !!{structure}~$\Struct{M}$ کجھ !!{sentence}s نوں سچ تے کجھ نوں
جھوٹا بناؤندی اے۔ اوہدے سچ بنائے سارے !!{sentence}s دے سیٹ نوں
اوہدا \emph{نظریہ} آکھیا جاندا اے۔ ایہ سیٹ واقعی اک نظریہ اے، کیوں جے
اوہدے توں جیہڑی وی گل منطقی طور اُتے نکلے، اوہ اوہدے سارے نمونیاں
وچ سچ ہونی چاہیدی اے، جنہاں وچ~$\Struct{M}$ وی شامل اے۔ ایس لئی
اوہ نتیجہ وی ایس سارے سچ جملیاں دے سیٹ وچ آ جاندا اے۔

\begin{defn}
  !!a{structure}~$\Struct M$ دتی ہووے۔ $\Struct{M}$ دا
  \emph{نظریہ}، اوہناں !!{sentence}s دا سیٹ $\Theory{M}$ اے
  جیہڑے $\Struct{M}$ وچ سچ نیں؛ یعنی $\Theory{M} =
  \Setabs{!A}{\Sat{M}{!A}}$۔
\end{defn}

اسیں ``نظریہ'' دا لفظ غیر رسمی طور اُتے اوہناں !!{sentence}s دے
سیٹاں لئی وی ورتدے آں جنہاں دی کوئی مِتھی ہوئی تعبیر ہووے، چاہے اوہ
استنتاج دے لحاظ نال بند ہوون یا نہ ہوون۔

\begin{prop}
ہر $\Struct{M}$ لئی، $\Theory{M}$ مکمل اے۔
\end{prop}

\begin{proof}
ہر !!{sentence}~$!A$ لئی یا $\Sat{M}{!A}$ یا $\Sat{M}{\lnot !A}$،
ایس لئی یا $!A \in \Theory{M}$ یا $\lnot !A \in \Theory{M}$۔
\end{proof}

\begin{prop}\ollabel{prop:equiv}
  جے $\Struct{N} \models !A$ ہر $!A \in \Theory{M}$ لئی، تاں
  $\Struct{M} \elemequiv \Struct{N}$۔ ایتھے دوہاں ساختاں دی زبان اکّو اے۔
\end{prop}

\begin{proof}
کیوں جے $\Sat{N}{!A}$ ہر $!A \in \Theory{M}$ لئی، ایس لئی
$\Theory{M} \subseteq \Theory{N}$۔ جے $\Sat{N}{!A}$، تاں
$\Sat/{N}{\lnot !A}$؛ ایس لئی $\lnot !A
\notin \Theory{M}$۔ جے نفی پہلے نظریے وچ ہوندی، تاں دوسری ساخت وچ وی
سچ ہونی سی، پر ایتھے اوہ سچ نہیں۔ کیوں جے $\Theory{M}$ مکمل اے، ایس لئی $!A \in
\Theory{M}$۔ سو $\Theory{N} \subseteq \Theory{M}$، تے سانوں
$\Struct{M} \elemequiv \Struct{N}$ ملدا اے۔
\end{proof}

\begin{rem}\ollabel{remark:R}
  $\Struct{R} = \langle\Real, <\rangle$ نوں ویکھو: ایہ اوہ
  !!{structure} اے جیہدا دائرہ حقیقی اعداد دا سیٹ $\Real$ اے،
  تے اوہ !!{language} وچ اے جیہدے وچ صرف اک دو-تھانواں والا
  !!{predicate} اے، جس دی تعبیر حقیقی اعداد اُتے $<$ تعلق اے۔
  ظاہر اے کہ $\Struct{R}$ !!{nonenumerable} اے؛ پر چونکہ
  $\Theory{R}$ واضح طور اُتے موافق اے، L\"owenheim--Skolem قضیے
  مطابق ایہدا اک !!a{enumerable} نمونہ اے، مثلاً $\Struct{S}$؛
  تے \olref{prop:equiv} توں $\Struct{R}
  \equiv \Struct{S}$۔ اک دا دائرہ قابلِ شمار اے تے دوجے دا نہیں،
  ایس لئی اوہناں وچ ہم شکلی لئی درکار اک بہ اک مطابقت نہیں ہو سکدی۔
  ہور ایہ کہ $\Struct{R}$ تے $\Struct{S}$
  ہم شکل نہیں، ایس لئی ایہ مثال دکھاؤندی اے کہ
  \olref[iso]{thm:isom} دا الٹا دعویٰ عام طور اُتے درست نہیں۔
\end{rem}

\end{document}
